\documentclass[11pt]{article}

\usepackage{fontspec}
\setmainfont{DejaVu Serif}
\setsansfont{DejaVu Sans}
\setmonofont{DejaVu Sans Mono}

\usepackage[russian]{babel}
\usepackage[a4paper,margin=2.4cm]{geometry}
\usepackage{amsmath,amssymb}
\usepackage{boxedminipage}
\usepackage{enumitem}
\usepackage{fancyvrb}
\usepackage{titlesec}
\usepackage{hyperref}

\titleformat{\section}{\large\bfseries}{\thesection.}{0.6em}{}
\setlist{itemsep=2pt,topsep=3pt}

\newcommand{\imp}[1]{\overline{#1}}

\title{\bfseries Нахождение всех ядерных импликант булевой функции\\
методом SAT (без построения таблицы истинности)}
\author{}
\date{}

\begin{document}
\maketitle
\vspace{-2.5em}

\section{Сведение проверки ядерности к задаче выполнимости}

Пучком двоичного набора $\alpha$ называется
\[
\Pi(f,\alpha)=\{\,K_j \mid K_j(\alpha)=1\,\}.
\]

По теореме~3.2 из лекций:
\[
K_i \in \bigcap T_f
\iff
\exists\,\alpha:\ \Pi(f,\alpha)=\{K_i\}.
\]

То есть $K_i$ является ядерной импликантой тогда и только тогда, когда существует входной набор, на котором $K_i$ равна~1, а все остальные простые импликанты равны~0.

Поэтому для каждой $K_i$ строим формулу
\[
\boxed{\;F_i \;=\; K_i \;\land\; \bigwedge_{\substack{j=1\\ j\ne i}}^{m} \neg K_j.\;}
\]

Тогда
\[
\boxed{\;K_i \in \bigcap T_f \iff F_i \text{ выполнима}.\;}
\]

Здесь «выполнима» означает: существует набор значений переменных, на котором формула принимает значение~1.

\section{Построение $F_i$ сразу в виде КНФ}

SAT-решатели обычно принимают формулу в конъюнктивной нормальной форме (КНФ): конъюнкцию дизъюнктов, каждый из которых есть дизъюнкция литералов.

Пусть
\[
K_j = l_{j1}\land\cdots\land l_{jr_j},
\]
где каждый литерал $l_{jt}$ — переменная или её отрицание. По закону де Моргана
\[
\neg K_j = \neg l_{j1}\vee\cdots\vee\neg l_{jr_j}.
\]

Поэтому построение $F_i$ таково:

\begin{itemize}
\item для каждого литерала $l$ из $K_i$ добавляем единичный дизъюнкт $(l)$ — это гарантирует $K_i=1$;
\item для каждого $j\ne i$ добавляем один дизъюнкт
\[
(\neg l_{j1}\vee\cdots\vee\neg l_{jr_j}),
\]
что гарантирует $K_j=0$.
\end{itemize}

Например, если проверяется
\[
K_i = x_1\,\imp{x_3},
\]
то добавляются единичные дизъюнкты
\[
(x_1),\qquad (\imp{x_3}).
\]

Если другой терм —
\[
K_j = x_1 x_2 \,\imp{x_4},
\]
то добавляется дизъюнкт
\[
(\imp{x_1}\vee\imp{x_2}\vee x_4).
\]

Полученная формула уже является КНФ: \textbf{никакой таблицы истинности строить не нужно, дополнительных переменных не требуется.}

\section{Поиск выполняющего набора: алгоритм DPLL}

Сначала выполняется \textbf{единичное распространение} (unit propagation): если дизъюнкт содержит единственный литерал, этот литерал обязан быть истинным. После присваивания переменной удаляются выполненные дизъюнкты, а из остальных удаляются ставшие ложными литералы; процесс повторяется.

Например, в
\[
(x_1)\land(\imp{x_1}\vee x_2)
\]
сначала выводится $x_1=1$; второй дизъюнкт превращается в $(x_2)$, откуда $x_2=1$.

В ходе распространения:

\begin{itemize}
\item если возник \emph{пустой дизъюнкт} (все его литералы ложны) — текущая ветвь противоречива;
\item если все дизъюнкты выполнены — выполняющий набор найден; неприсвоенные переменные могут получить произвольные значения;
\item если остались невыполненные дизъюнкты, но единичных нет — выбирается неприсвоенная переменная, и рекурсивно проверяются оба значения $0$ и $1$.
\end{itemize}

Псевдокод:

\begin{Verbatim}[fontsize=\small,frame=single,framesep=6pt]
DPLL(множество дизъюнктов C, текущее частичное присваивание a):

    выполнить единичное распространение в C, попутно обновляя a

    если C содержит пустой дизъюнкт:
        вернуть UNSAT

    если C пусто:
        произвольно доопределить неприсвоенные переменные в a
        вернуть SAT и набор a

    выбрать неприсвоенную переменную x из C

    R := DPLL(C ∪ {(x)}, a)
    если R = SAT:
        вернуть R

    вернуть DPLL(C ∪ {(¬x)}, a)
\end{Verbatim}

Процесс работает непосредственно с дизъюнктами и частичными присваиваниями. При конфликте отсекается целая ветвь поиска, соответствующая данному частичному присваиванию, — перебирать все входные наборы под ней не требуется.

\section{Алгоритм нахождения всех ядерных импликант}

Литералы кодируются целыми числами со знаком: $x_k$ — числом $k$, $\imp{x_k}$ — числом $-k$. Например,
\[
x_1\imp{x_3}x_5 \quad\longleftrightarrow\quad [1,-3,5],
\]
отрицание литерала — смена знака числа.

Полный алгоритм:

\begin{Verbatim}[fontsize=\small,frame=single,framesep=6pt]
Вход:  все простые импликанты K_1, ..., K_m
Выход: множество ядерных импликант Core

Core := ∅

для i = 1, ..., m:

    C := пустое множество дизъюнктов

    для каждого литерала l из K_i:
        добавить в C единичный дизъюнкт [l]

    для каждого j ≠ i:
        добавить в C дизъюнкт [-l : l ∈ K_j]

    вызвать SAT-решатель для C

    если результат SAT:
        добавить K_i в Core
        сохранить возвращённый набор α_i

    если результат UNSAT:
        K_i в Core не включать

вернуть Core
\end{Verbatim}

Сохранённый набор $\alpha_i$ — свидетель ядерности $K_i$: по нему можно непосредственно проверить
\[
K_i(\alpha_i)=1,\qquad K_j(\alpha_i)=0\quad (j\ne i).
\]

При проверке каждого терма всегда используются \emph{все} остальные термы исходной сокращённой ДНФ: уже признанные неядерными термы из последующих запросов удалять нельзя.

\section{Пример}

Рассмотрим сокращённую ДНФ
\[
A_f = x_1x_2 \vee \imp{x_1}x_3 \vee x_2x_3
\]
со всеми простыми импликантами
\[
K_1=x_1x_2,\qquad K_2=\imp{x_1}x_3,\qquad K_3=x_2x_3.
\]

\textbf{Проверка $K_1$.} Строим
\[
F_1 =
x_1\land x_2
\land(x_1\vee\imp{x_3})
\land(\imp{x_2}\vee\imp{x_3}).
\]
Единичные дизъюнкты дают $x_1=1,\ x_2=1$; тогда последний дизъюнкт требует $x_3=0$. Выполняющий набор:
\[
\alpha_1=(1,1,0).
\]
Значит, $K_1$ — ядерная импликанта.

\textbf{Проверка $K_2$.} Строим
\[
F_2 =
\imp{x_1}\land x_3
\land(\imp{x_1}\vee\imp{x_2})
\land(\imp{x_2}\vee\imp{x_3}).
\]
Единичное распространение даёт
\[
x_1=0,\qquad x_3=1,\qquad x_2=0.
\]
Набор $\alpha_2=(0,0,1)$ выполняет $F_2$, поэтому $K_2$ тоже ядерная.

\textbf{Проверка $K_3$.} Строим
\[
F_3 =
x_2\land x_3
\land(\imp{x_1}\vee\imp{x_2})
\land(x_1\vee\imp{x_3}).
\]
Из $x_2=1,\ x_3=1$ два последних дизъюнкта превращаются соответственно в
\[
\imp{x_1},\qquad x_1,
\]
что противоречиво. Следовательно, $F_3$ невыполнима, и $K_3$ не является ядерной.

Итог:
\[
\boxed{\;\bigcap T_f = \{\,x_1x_2,\ \imp{x_1}x_3\,\}.\;}
\]

\section{Доказательство корректности}

Если алгоритм включает $K_i$ в \texttt{Core}, то решатель нашёл выполняющий набор $\alpha$ формулы $F_i$. По построению $F_i$:
\[
K_i(\alpha)=1,\qquad K_j(\alpha)=0\quad (j\ne i).
\]
Так как вход содержит \emph{все} простые импликанты, получаем
\[
\Pi(f,\alpha)=\{K_i\},
\]
то есть $K_i$ — ядерная импликанта, $K_i\in\bigcap T_f$.

Обратно, если $K_i\in\bigcap T_f$, то по теореме из лекций существует набор $\alpha$ с
\[
\Pi(f,\alpha)=\{K_i\}.
\]
Этот набор выполняет $F_i$, поэтому полная SAT-процедура обязана признать $F_i$ выполнимой и включить $K_i$ в результат.

Следовательно, алгоритм возвращает в точности
\[
\boxed{\;\operatorname{Core} = \bigcap T_f.\;}
\]

\section{Оценка объёма вычислений}

Пусть $K_i$ содержит $r_i$ литералов ($r_i\le n$). Каждый запрос содержит:

\begin{itemize}
\item не более $n=100$ переменных;
\item $r_i$ единичных дизъюнктов;
\item $m-1=999$ дизъюнктов — отрицаний остальных термов.
\end{itemize}

Итого не более
\[
100+999=1099
\]
дизъюнктов на запрос; всего требуется $m=1000$ SAT-запросов.

Если суммарное число вхождений литералов равно
\[
L=\sum_{j=1}^{m} r_j,
\]
то каждый запрос содержит ровно $L$ вхождений литералов, построение одного запроса — $O(L)$, всех запросов — $O(mL)$. Время решения зависит от сложности конкретных SAT-экземпляров.

Метод не строит таблицу истинности из $2^{100}$ строк, а проверяет существование «единственного» набора посредством единичного распространения, ветвления и обработки конфликтов.

\end{document}
